\chapter{Prealbumin (Transthyretin)}

\section*{What Is This Test?}
Prealbumin (transthyretin, TTR) is a small (55 kDa) transport protein synthesized exclusively by the liver. It transports thyroxine and retinol (vitamin A) and serves as a sensitive marker of visceral protein synthesis. With a half-life of approximately 2 days --- far shorter than albumin (20 days) --- prealbumin reflects acute changes in nutritional status and hepatic protein synthesis. It is widely used to detect malnutrition, monitor nutritional repletion in hospitalized patients, and assess liver synthetic function. Prealbumin may also serve as a marker of inflammation (negative acute-phase reactant) and is the precursor protein in hereditary and wild-type transthyretin amyloidosis.

\section*{Alternative Names}
Prealbumin; Transthyretin; TTR; Thyroxine-Binding Prealbumin; TBPA

\section*{Principle}
Immunonephelometric or immunoturbidimetric assay using antibody specific for prealbumin. Antigen-antibody complexes scatter light proportionally to prealbumin concentration. Serum separator or EDTA tube is acceptable; fasting not required.

\section*{History}
Prealbumin was named for its electrophoretic mobility just ahead of albumin. It was identified as a carrier protein for thyroxine and retinol in the 1960s. The hepatology and nutrition communities adopted its clinical use as a short-half-life nutritional marker in the 1980s. The hereditary amyloidosis community recognized transthyterin mutations as causative in ATTR amyloidosis in the 1990s.

\section*{Why This Test Is Ordered}
\begin{itemize}
  \item Nutritional assessment in hospitalized patients (especially before GI surgery, dialysis, ICU stay)
  \item Monitoring response to enteral or parenteral nutrition
  \item Viseral protein synthesis assessment in liver disease
  \item Identification of protein-energy malnutrition in chronic illnesses, dialysis, elderly
  \item Investigation of older adults with weight loss and frailty
  \item Cardiology workup for patients with restrictive cardiomyopathy (assessment of transthyretin amyloidosis)
\end{itemize}

\section*{Is This a Primary or Secondary Test}
Secondary test for nutritional and liver function assessment; complementary to albumin.

\section*{How This Test Is Performed}
Venous blood drawn in serum separator tube; results available within 4--24 hours. Serial measurements every 3--7 days are useful for monitoring nutrition repletion.

\section*{What Preparation Is Needed}
None. Fasting is not required.

\section*{Contraindications and Precautions}
Prealbumin is a negative acute-phase reactant: inflammation reduces values as an unrelated consequence of cytokine-mediated hepatic reprioritization; concomitant C-reactive protein (CRP) measurement helps interpret. Renal failure increases levels due to reduced clearance; pregnancy and estrogen therapy raise values.

\section*{Risks}
None.

\section*{What Happens After the Test}
Values $<$ 15 mg/dL suggest malnutrition; values $<$ 10 mg/dL severe malnutrition. Increasing prealbumin over 7--14 days of nutrition support confirms adequate repletion. Interpretation flagged for inflammation (CRP elevated).

\section*{Understanding Results}

\subsection*{Below Normal}
\begin{itemize}
  \item Mild malnutrition: 10--15 mg/dL
  \item Moderate malnutrition: 5--10 mg/dL
  \item Severe malnutrition: $<$ 5 mg/dL
  \item Liver disease: decreased synthesis (cirrhosis, hepatitis, liver failure)
  \item Acute inflammation: protein redirects to acute-phase reactants (ferritin, CRP)
  \item Hyperthyroidism: increased turnover
  \item Nephrotic syndrome or protein-losing enteropathy
\end{itemize}

\subsection*{Above Normal}
\begin{itemize}
  \item Renal failure (decreased clearance): chronic kidney disease
  \item High-dose corticosteroid therapy
  \item Pregnancy and estrogen therapy: increased hepatic synthesis
  \item Hodgkin lymphoma or HPT hormone use
\end{itemize}

\section*{Normal Results}
\begin{itemize}
  \item \textbf{Adult (serum):} 15--35 mg/dL (150--350 mg/L)
  \item \textbf{Elderly:} slightly lower (15--25 mg/dL)
  \item \textbf{Children:} 12--30 mg/dL
  \item \textbf{Pregnancy (3rd trimester):} typically elevated
  \item \textbf{Nutritional repletion goal:} Rise of $\geq$ 2 mg/dL within 7--14 days suggests adequate nutrition support
\end{itemize}

\section*{Follow-up Tests}
\begin{itemize}
  \item \textbf{Albumin and total protein:} Long-term protein synthesis markers
  \item \textbf{C-reactive protein (CRP):} Confounds acute-phase suppression
  \item \textbf{Comprehensive nutritional assessment:} Anthropometry, dietary intake, micronutrient levels
  \item \textbf{ Liver function tests (ALT, AST, INR):} If hepatic synthetic failure suspected
  \item \textbf{Renal function (creatinine, BUN, urinalysis):} High prealbumin suggests renal failure
  \item \textbf{Transthyretin genetic (TTR) testing:} If hereditary amyloidosis suspected (clinical: cardiomyopathy, neuropathy)
  \item \textbf{Pyrophosphate/C-PYP scintigraphy (PYP scan):} For transthyretin amyloid cardiomyopathy diagnosis
\end{itemize}

\section*{References}
\begin{enumerate}
  \item Bernstein L. Prealbumin as a marker of nutritional status. Nutr Clin Pract. 2012;27(5):681-684.
  \item Bernstein LH, et al. Prealbumin: A critical assay discussion. Nutr Clin Care. 2005;8(3):129-136.
  \item Conceição I, et al. Red-flag index for ATTR amyloidosis. Eur J Neurol. 2017;24(6):e43.
  \item Schreiber SS, et al. Transthyretin: from discovery to amyloid disease therapy. Front Med. 2022;9:856139.
\end{enumerate}

\clearpage